% 補足資料の日本語版．英語版 ../supplementary.tex と節・図・表・式・数値を 1 対 1 に対応させる．
\documentclass[paper=a4,fontsize=9.5pt,twocolumn,line_length=26zw,number_of_lines=50,head_space=17mm,baselineskip=1.55zw,column_gap=2zw]{jlreq}
\usepackage[ipaex]{luatexja-preset}
\usepackage{amsmath,amssymb}
\usepackage{graphicx}
\usepackage{booktabs}
\usepackage[nocompress]{cite}
\usepackage{url}
\usepackage[hidelinks]{hyperref}

% \emph の中の和文はゴシックにする（欧文は斜体のまま）．
\let\jaorigemph\emph
\renewcommand{\emph}[1]{\jaorigemph{\gtfamily #1}}

% 自動生成．手で編集しない．
% 生成元: results/numbers.json ／ 生成器: scripts/gen_numbers_tex.py
% 本文の実験由来の数値は \NUM{key} で書く．未登録の key はビルド時にエラーになる．
\makeatletter
\providecommand{\NUM}[1]{%
  \@ifundefined{NUM@#1}{%
    \GenericError{}{numbers.tex: key '#1' is not registered in results/numbers.json}{}{}%
    \textbf{??}%
  }{\csname NUM@#1\endcsname}%
}
\makeatother

\newcommand{\FIG}[2]{\IfFileExists{../figures/#1.pdf}{\includegraphics[width=#2]{../figures/#1.pdf}}{\fbox{\parbox[c][6em][c]{\dimexpr#2-2\fboxsep-2\fboxrule\relax}{\centering\footnotesize figure \texttt{#1} (pending)}}}}
\renewcommand{\thesection}{S\arabic{section}}
\renewcommand{\thetable}{S\arabic{table}}
\renewcommand{\thefigure}{S\arabic{figure}}
\renewcommand{\theequation}{S\arabic{equation}}

\title{「少数のターゲットラベルが最初に直すもの：\\コーパス横断のマルチモーダル感情認識における\\判断基準と向きの分析，\\およびモダリティをまたぐ弱いラベルの転移」の補足資料}
\author{著者名（未記入）}
\date{}

% 段の下端をそろえるために式の前後が引き伸ばされるのを防ぐ．
\raggedbottom
% 図と表の説明文：ラベル（図1，表1）はゴシックの太字，本文は明朝で small．
\jlreqsetup{caption_font={\small},caption_label_font={\small\sffamily\gtfamily\bfseries}}

\begin{document}
\maketitle

\section{特徴の抽出}
\label{sec:supp-features}

\emph{顔．}
各動画から$\lfloor 5\,\mathrm{fps}/25 \rfloor$フレームに1枚を取り出し，これは毎秒5から6フレームになる．
発話あたり最大$400$フレームを等間隔に抜き出す．
IEMOCAPの動画は2人の話者を左右に並べて映しているので，発話の話者を映す半分を残す．
YOLOv8n-face\cite{yolov8face}が確信度のしきい値$0.25$で顔を検出し，各フレームの最大の顔を残して$224\times224$画素の大きさにそろえる．
Ryuminaら\cite{ryumina2022facial}のResNet-50が各切り出しを$512$次元に写し，第2のエンコーダであるCLIP ViT-L/14\cite{radford2021clip}の画像エンコーダが同じ切り出し（枠として保存したもの）を$1024$次元に写す．
フレームの埋め込みを，$10$フレームずつずらした$20$顔フレームの区間の中で平均し，$10$フレームより短い最後の区間は捨て，区間の平均をさらに平均する．
区間の間の標準偏差は使わない．
顔のフレームが$20$未満の発話は区間が1つで標準偏差が0になり，この特徴が発話の長さを表してしまうからである．

\emph{音声．}
波形を$16$~kHzのモノラルに再標本化し，$60$~sで打ち切り，平均0，分散1に正規化して，WavLM-Large\cite{chen2022wavlm}またはHuBERT-Large\cite{hsu2021hubert}に通す．
最終層を，パディングを含めずに時間で平均する（$1024$次元）．

\emph{テキスト．}
書き起こしの文字コードの誤りを直し，$256$トークンで打ち切り，RoBERTa-base\cite{liu2019roberta}またはDeBERTa-v3-base\cite{he2023debertav3}に通す．
最終層をトークンで平均する（$768$次元）．

\emph{発話数．}
除外の後，音声の特徴は\NUM{data-n-total-wavlm}発話，顔の特徴は\NUM{data-n-total-resnet}発話，テキストの特徴は\NUM{data-n-total-roberta}発話を含む．
音声と顔を融合すると，IEMOCAPの即興のセッションと台本のセッション，CMU-MOSEI，MELD，AFEW，CREMA-Dでそれぞれ\NUM{data-n-impro-av}，\NUM{data-n-script-av}，\NUM{data-n-mosei-av}，\NUM{data-n-meld-av}，\NUM{data-n-afew-av}，\NUM{data-n-cremad-av}発話が残る．
3つのモダリティすべてを融合すると，テキストのある4つの群でそれぞれ\NUM{data-n-impro-avt}，\NUM{data-n-script-avt}，\NUM{data-n-mosei-avt}，\NUM{data-n-meld-avt}発話が残る．

\section{導出}
\label{sec:supp-derivations}

\emph{判断基準のコスト．}
対ごとの判定の向きを固定し，スコア差$z$が平均$\mu_k$，$\mu_l$と共通の標準偏差$\sigma$をもつガウス分布に従うとする．
2つの感情を釣り合わせると，判断基準は$(\mu_k + \mu_l)/2$に置かれる．
感情ごとに$K$本のラベルから推定すると，2つのクラス平均の分散はそれぞれ$\sigma^2/K$なので，中点の分散は$\sigma^2/(2K)$となり，残るずれは$\mathbb{E}[o^2] = 1/(2K)$を満たす．
$o=0$で$\partial a/\partial o = 0$かつ$\partial^2 a/\partial o^2 = -(d/2)\,\phi(d/2)$なので，
\begin{equation}
a(d, 0) - a(d, o) \approx \frac{d}{4}\,\phi\!\left(\frac{d}{2}\right) o^2
\end{equation}
が成り立ち，期待値をとると本文の判断基準のコストが得られる．
係数$(d/4)\phi(d/2)$は$d=2$で最大となり，そのとき$\phi(1)/2$に等しい．

\emph{向きのコスト．}
感情ごとに$K$本のラベルから推定した線形の規則では，推定した向きと最良の向きのなす角$\theta$は，1次の近似で$\tan^2\theta \approx q/K$を満たす．
ここで$q$は実効次元である．
$p$個の特徴と2乗マハラノビス距離$\Delta^2$をもつ標本線形判別では，誤分類率の展開の1次の項から$q = 2(p-1)/\Delta^2$が得られる\cite{okamoto1963asymptotic,raudys1998intrinsic}．
推定した向きに沿った識別力は$\Delta\cos\theta = \Delta/\sqrt{1 + \tan^2\theta}$であり，これから$d(K) = \Delta/\sqrt{1 + q/K}$が得られる．
$\ell_2$罰則があると実効次元は$p$より小さく，本論文では$q$をターゲットのみの学習曲線に当てはめる．

\emph{ソースの価値．}
向きの推定の雑音と信号の比を$\tan^2\theta$と書く．
$K$本のラベルからのターゲットの推定の比は$q/K$であり，効率$e = \cos\theta_s$のソースの向きの比は$(1-e^2)/e^2$である．
2つの独立な推定を分散の逆数で重みづけして組み合わせると，比は$\left[K/q + e^2/(1-e^2)\right]^{-1}$となる．
これは，$n = q e^2/(1-e^2)$としたときの$K + n$本のラベルからのターゲットのみの推定の比に等しい．

\emph{合併による判断基準の修正．}
合併の目的関数を，ソースのみの目的関数に$N_t$個のターゲットの点の損失を加えたものとして書き，ソースのみの解からニュートン法の1ステップをとる．
損失の曲率をドメインごとの平均$\kappa$で近似すると，このステップはターゲットの平均でのスコアを，ソースのみのずれの$\rho/(1+\rho)$倍だけ変える．
ここで$\rho$は本文の式で与えられ，項$1/(N_s\kappa_s)$は切片から，項$\boldsymbol{\delta}^\top(A_s + I/C)^{-1}\boldsymbol{\delta}$は重みから生じる．
この近似は，ターゲットの点の損失をずれについて2次まで残し，向きをソースから取り，ターゲットの点そのものの散らばりを省く．
$128$次元のガウス分布の特徴によるシミュレーションは，この式を確かめる（図\ref{fig:supp-lambda}に実際のコーパスの対での確認を示す）．

\section{第I部の追加の結果}
\label{sec:supp-part1}

\begin{figure}[t]
\centering
\FIG{lambda}{\columnwidth}
\caption{合併の分類器の判断基準の修正$\lambda$の予言値と測定値．順序対と$K \in \{1, 5, 10\}$ごとに示す．}
\label{fig:supp-lambda}
\end{figure}

\begin{figure}[t]
\centering
\FIG{centering}{\columnwidth}
\caption{中心化の有無によるソースのみの分類器，感情ごとに$K=5$本のラベルで切片だけを当てはめ直したもの，同じラベルによる合併を，WavLMの音声，HuBERTの音声，顔について示す．}
\label{fig:supp-centering}
\end{figure}

\section{比較の表の全体}
\label{sec:supp-tables}

% 英語版 ../../supp/tables.tex（scripts/make_supp_tables.py の生成物）の日本語版．行の \NUM は英語版と同じ．
% 英語版の表が作り直されたら，説明文と見出しのほかはその行を写し直す．
\begin{table*}[t]
\caption{単一モダリティの生徒における，調整版の比較手法と提案法（提案法と提案法の拡張版）のUAR（\%）．ターゲットと10回の反復にわたって平均した．ST：自己訓練，HPL：ハードな擬似ラベル，EM：エントロピー最小化，KD：温度つきの蒸留，GD：一般化蒸留，CoT：共訓練，CoR：共正則化，PA：EMによる事前確率の調整，RB：出力の再均衡化．}
\label{tab:supp-mod}
\centering\footnotesize\setlength{\tabcolsep}{3pt}
\begin{tabular}{@{}llrrrrrrrrrrrr@{}}
\toprule
構成 & 生徒，$K$ & 出発点 & ST & HPL & EM & KD & GD & CoT & CoR & PA & RB & 提案法 & 提案法の拡張版 \\
\midrule
AVT & A，0 & \NUM{tab-avt-a-k0-e39} & \NUM{tab-avt-a-k0-self} & \NUM{tab-avt-a-k0-hard} & \NUM{tab-avt-a-k0-entmin} & \NUM{tab-avt-a-k0-kd} & \NUM{tab-avt-a-k0-gd} & \NUM{tab-avt-a-k0-cotrain} & \NUM{tab-avt-a-k0-coreg} & \NUM{tab-avt-a-k0-saerens} & \NUM{tab-avt-a-k0-massbal} & \NUM{tab-avt-a-k0-ours} & \NUM{tab-avt-a-k0-oursx} \\
 & A，5 & \NUM{tab-avt-a-k5-e39} & \NUM{tab-avt-a-k5-self} & \NUM{tab-avt-a-k5-hard} & \NUM{tab-avt-a-k5-entmin} & \NUM{tab-avt-a-k5-kd} & \NUM{tab-avt-a-k5-gd} & \NUM{tab-avt-a-k5-cotrain} & \NUM{tab-avt-a-k5-coreg} & \NUM{tab-avt-a-k5-saerens} & \NUM{tab-avt-a-k5-massbal} & \NUM{tab-avt-a-k5-ours} & \NUM{tab-avt-a-k5-oursx} \\
AVT & V，0 & \NUM{tab-avt-v-k0-e39} & \NUM{tab-avt-v-k0-self} & \NUM{tab-avt-v-k0-hard} & \NUM{tab-avt-v-k0-entmin} & \NUM{tab-avt-v-k0-kd} & \NUM{tab-avt-v-k0-gd} & \NUM{tab-avt-v-k0-cotrain} & \NUM{tab-avt-v-k0-coreg} & \NUM{tab-avt-v-k0-saerens} & \NUM{tab-avt-v-k0-massbal} & \NUM{tab-avt-v-k0-ours} & \NUM{tab-avt-v-k0-oursx} \\
 & V，5 & \NUM{tab-avt-v-k5-e39} & \NUM{tab-avt-v-k5-self} & \NUM{tab-avt-v-k5-hard} & \NUM{tab-avt-v-k5-entmin} & \NUM{tab-avt-v-k5-kd} & \NUM{tab-avt-v-k5-gd} & \NUM{tab-avt-v-k5-cotrain} & \NUM{tab-avt-v-k5-coreg} & \NUM{tab-avt-v-k5-saerens} & \NUM{tab-avt-v-k5-massbal} & \NUM{tab-avt-v-k5-ours} & \NUM{tab-avt-v-k5-oursx} \\
AVT & T，0 & \NUM{tab-avt-t-k0-e39} & \NUM{tab-avt-t-k0-self} & \NUM{tab-avt-t-k0-hard} & \NUM{tab-avt-t-k0-entmin} & \NUM{tab-avt-t-k0-kd} & \NUM{tab-avt-t-k0-gd} & \NUM{tab-avt-t-k0-cotrain} & \NUM{tab-avt-t-k0-coreg} & \NUM{tab-avt-t-k0-saerens} & \NUM{tab-avt-t-k0-massbal} & \NUM{tab-avt-t-k0-ours} & \NUM{tab-avt-t-k0-oursx} \\
 & T，5 & \NUM{tab-avt-t-k5-e39} & \NUM{tab-avt-t-k5-self} & \NUM{tab-avt-t-k5-hard} & \NUM{tab-avt-t-k5-entmin} & \NUM{tab-avt-t-k5-kd} & \NUM{tab-avt-t-k5-gd} & \NUM{tab-avt-t-k5-cotrain} & \NUM{tab-avt-t-k5-coreg} & \NUM{tab-avt-t-k5-saerens} & \NUM{tab-avt-t-k5-massbal} & \NUM{tab-avt-t-k5-ours} & \NUM{tab-avt-t-k5-oursx} \\
AVT（第2） & A，0 & \NUM{tab-w2-a-k0-e39} & \NUM{tab-w2-a-k0-self} & \NUM{tab-w2-a-k0-hard} & \NUM{tab-w2-a-k0-entmin} & \NUM{tab-w2-a-k0-kd} & \NUM{tab-w2-a-k0-gd} & \NUM{tab-w2-a-k0-cotrain} & \NUM{tab-w2-a-k0-coreg} & \NUM{tab-w2-a-k0-saerens} & \NUM{tab-w2-a-k0-massbal} & \NUM{tab-w2-a-k0-ours} & \NUM{tab-w2-a-k0-oursx} \\
 & A，5 & \NUM{tab-w2-a-k5-e39} & \NUM{tab-w2-a-k5-self} & \NUM{tab-w2-a-k5-hard} & \NUM{tab-w2-a-k5-entmin} & \NUM{tab-w2-a-k5-kd} & \NUM{tab-w2-a-k5-gd} & \NUM{tab-w2-a-k5-cotrain} & \NUM{tab-w2-a-k5-coreg} & \NUM{tab-w2-a-k5-saerens} & \NUM{tab-w2-a-k5-massbal} & \NUM{tab-w2-a-k5-ours} & \NUM{tab-w2-a-k5-oursx} \\
AVT（第2） & V，0 & \NUM{tab-w2-v-k0-e39} & \NUM{tab-w2-v-k0-self} & \NUM{tab-w2-v-k0-hard} & \NUM{tab-w2-v-k0-entmin} & \NUM{tab-w2-v-k0-kd} & \NUM{tab-w2-v-k0-gd} & \NUM{tab-w2-v-k0-cotrain} & \NUM{tab-w2-v-k0-coreg} & \NUM{tab-w2-v-k0-saerens} & \NUM{tab-w2-v-k0-massbal} & \NUM{tab-w2-v-k0-ours} & \NUM{tab-w2-v-k0-oursx} \\
 & V，5 & \NUM{tab-w2-v-k5-e39} & \NUM{tab-w2-v-k5-self} & \NUM{tab-w2-v-k5-hard} & \NUM{tab-w2-v-k5-entmin} & \NUM{tab-w2-v-k5-kd} & \NUM{tab-w2-v-k5-gd} & \NUM{tab-w2-v-k5-cotrain} & \NUM{tab-w2-v-k5-coreg} & \NUM{tab-w2-v-k5-saerens} & \NUM{tab-w2-v-k5-massbal} & \NUM{tab-w2-v-k5-ours} & \NUM{tab-w2-v-k5-oursx} \\
AVT（第2） & T，0 & \NUM{tab-w2-t-k0-e39} & \NUM{tab-w2-t-k0-self} & \NUM{tab-w2-t-k0-hard} & \NUM{tab-w2-t-k0-entmin} & \NUM{tab-w2-t-k0-kd} & \NUM{tab-w2-t-k0-gd} & \NUM{tab-w2-t-k0-cotrain} & \NUM{tab-w2-t-k0-coreg} & \NUM{tab-w2-t-k0-saerens} & \NUM{tab-w2-t-k0-massbal} & \NUM{tab-w2-t-k0-ours} & \NUM{tab-w2-t-k0-oursx} \\
 & T，5 & \NUM{tab-w2-t-k5-e39} & \NUM{tab-w2-t-k5-self} & \NUM{tab-w2-t-k5-hard} & \NUM{tab-w2-t-k5-entmin} & \NUM{tab-w2-t-k5-kd} & \NUM{tab-w2-t-k5-gd} & \NUM{tab-w2-t-k5-cotrain} & \NUM{tab-w2-t-k5-coreg} & \NUM{tab-w2-t-k5-saerens} & \NUM{tab-w2-t-k5-massbal} & \NUM{tab-w2-t-k5-ours} & \NUM{tab-w2-t-k5-oursx} \\
\bottomrule
\end{tabular}
\end{table*}

\begin{table*}[t]
\caption{融合の生徒における，調整版の比較手法と提案法（提案法と提案法の拡張版）のUAR（\%）．CoT-FとCoT-P：融合のビューまたは積のビューによる共訓練．その他の略号は表\ref{tab:supp-mod}と同じ．}
\label{tab:supp-fused}
\centering\footnotesize\setlength{\tabcolsep}{3pt}
\begin{tabular}{@{}llrrrrrrrrrrr@{}}
\toprule
構成 & $K$ & 出発点 & ST & EM & KD & CoT-F & CoT-P & CoR & PA & RB & 提案法 & 提案法の拡張版 \\
\midrule
AVT & 0 & \NUM{tab-avt-f-k0-e39} & \NUM{tab-avt-f-k0-self} & \NUM{tab-avt-f-k0-entmin} & \NUM{tab-avt-f-k0-kd} & \NUM{tab-avt-f-k0-cotrain-fused} & \NUM{tab-avt-f-k0-cotrain-prod} & \NUM{tab-avt-f-k0-coreg-sum} & \NUM{tab-avt-f-k0-saerens} & \NUM{tab-avt-f-k0-massbal} & \NUM{tab-avt-f-k0-ours} & \NUM{tab-avt-f-k0-oursx} \\
 & 5 & \NUM{tab-avt-f-k5-e39} & \NUM{tab-avt-f-k5-self} & \NUM{tab-avt-f-k5-entmin} & \NUM{tab-avt-f-k5-kd} & \NUM{tab-avt-f-k5-cotrain-fused} & \NUM{tab-avt-f-k5-cotrain-prod} & \NUM{tab-avt-f-k5-coreg-sum} & \NUM{tab-avt-f-k5-saerens} & \NUM{tab-avt-f-k5-massbal} & \NUM{tab-avt-f-k5-ours} & \NUM{tab-avt-f-k5-oursx} \\
AV & 0 & \NUM{tab-av-f-k0-e39} & \NUM{tab-av-f-k0-self} & \NUM{tab-av-f-k0-entmin} & \NUM{tab-av-f-k0-kd} & \NUM{tab-av-f-k0-cotrain-fused} & \NUM{tab-av-f-k0-cotrain-prod} & \NUM{tab-av-f-k0-coreg-sum} & \NUM{tab-av-f-k0-saerens} & \NUM{tab-av-f-k0-massbal} & \NUM{tab-av-f-k0-ours} & \NUM{tab-av-f-k0-oursx} \\
 & 5 & \NUM{tab-av-f-k5-e39} & \NUM{tab-av-f-k5-self} & \NUM{tab-av-f-k5-entmin} & \NUM{tab-av-f-k5-kd} & \NUM{tab-av-f-k5-cotrain-fused} & \NUM{tab-av-f-k5-cotrain-prod} & \NUM{tab-av-f-k5-coreg-sum} & \NUM{tab-av-f-k5-saerens} & \NUM{tab-av-f-k5-massbal} & \NUM{tab-av-f-k5-ours} & \NUM{tab-av-f-k5-oursx} \\
VT & 0 & \NUM{tab-vt-f-k0-e39} & \NUM{tab-vt-f-k0-self} & \NUM{tab-vt-f-k0-entmin} & \NUM{tab-vt-f-k0-kd} & \NUM{tab-vt-f-k0-cotrain-fused} & \NUM{tab-vt-f-k0-cotrain-prod} & \NUM{tab-vt-f-k0-coreg-sum} & \NUM{tab-vt-f-k0-saerens} & \NUM{tab-vt-f-k0-massbal} & \NUM{tab-vt-f-k0-ours} & \NUM{tab-vt-f-k0-oursx} \\
 & 5 & \NUM{tab-vt-f-k5-e39} & \NUM{tab-vt-f-k5-self} & \NUM{tab-vt-f-k5-entmin} & \NUM{tab-vt-f-k5-kd} & \NUM{tab-vt-f-k5-cotrain-fused} & \NUM{tab-vt-f-k5-cotrain-prod} & \NUM{tab-vt-f-k5-coreg-sum} & \NUM{tab-vt-f-k5-saerens} & \NUM{tab-vt-f-k5-massbal} & \NUM{tab-vt-f-k5-ours} & \NUM{tab-vt-f-k5-oursx} \\
AT & 0 & \NUM{tab-at-f-k0-e39} & \NUM{tab-at-f-k0-self} & \NUM{tab-at-f-k0-entmin} & \NUM{tab-at-f-k0-kd} & \NUM{tab-at-f-k0-cotrain-fused} & \NUM{tab-at-f-k0-cotrain-prod} & \NUM{tab-at-f-k0-coreg-sum} & \NUM{tab-at-f-k0-saerens} & \NUM{tab-at-f-k0-massbal} & \NUM{tab-at-f-k0-ours} & \NUM{tab-at-f-k0-oursx} \\
 & 5 & \NUM{tab-at-f-k5-e39} & \NUM{tab-at-f-k5-self} & \NUM{tab-at-f-k5-entmin} & \NUM{tab-at-f-k5-kd} & \NUM{tab-at-f-k5-cotrain-fused} & \NUM{tab-at-f-k5-cotrain-prod} & \NUM{tab-at-f-k5-coreg-sum} & \NUM{tab-at-f-k5-saerens} & \NUM{tab-at-f-k5-massbal} & \NUM{tab-at-f-k5-ours} & \NUM{tab-at-f-k5-oursx} \\
AVT（第2） & 0 & \NUM{tab-w2-f-k0-e39} & \NUM{tab-w2-f-k0-self} & \NUM{tab-w2-f-k0-entmin} & \NUM{tab-w2-f-k0-kd} & \NUM{tab-w2-f-k0-cotrain-fused} & \NUM{tab-w2-f-k0-cotrain-prod} & \NUM{tab-w2-f-k0-coreg-sum} & \NUM{tab-w2-f-k0-saerens} & \NUM{tab-w2-f-k0-massbal} & \NUM{tab-w2-f-k0-ours} & \NUM{tab-w2-f-k0-oursx} \\
 & 5 & \NUM{tab-w2-f-k5-e39} & \NUM{tab-w2-f-k5-self} & \NUM{tab-w2-f-k5-entmin} & \NUM{tab-w2-f-k5-kd} & \NUM{tab-w2-f-k5-cotrain-fused} & \NUM{tab-w2-f-k5-cotrain-prod} & \NUM{tab-w2-f-k5-coreg-sum} & \NUM{tab-w2-f-k5-saerens} & \NUM{tab-w2-f-k5-massbal} & \NUM{tab-w2-f-k5-ours} & \NUM{tab-w2-f-k5-oursx} \\
\bottomrule
\end{tabular}
\end{table*}

\begin{table}[t]
\caption{ターゲットと反復の組にわたる，提案法の拡張版から調整版の最も強い比較手法を引いた差．平均（ポイント），効果量$d_z$，Wilcoxonの符号順位検定のHolmの補正後の$p$，平均が正のターゲットの数（正）．}
\label{tab:supp-stats}
\centering\footnotesize\setlength{\tabcolsep}{2.5pt}
\begin{tabular}{@{}llrrrrrrrr@{}}
\toprule
 & & \multicolumn{4}{c}{$K=0$} & \multicolumn{4}{c}{$K=5$} \\ \cmidrule(lr){3-6}\cmidrule(l){7-10}
構成 & 生徒 & 平均 & $d_z$ & $p$ & 正 & 平均 & $d_z$ & $p$ & 正 \\
\midrule
AVT & A & \NUM{st-oursx-best-avt-a-k0-mean} & \NUM{st-oursx-best-avt-a-k0-dz} & $<$0.001 & \NUM{st-oursx-best-avt-a-k0-tgtpos} & \NUM{st-oursx-best-avt-a-k5-mean} & \NUM{st-oursx-best-avt-a-k5-dz} & $<$0.001 & \NUM{st-oursx-best-avt-a-k5-tgtpos} \\
AVT & V & \NUM{st-oursx-best-avt-v-k0-mean} & \NUM{st-oursx-best-avt-v-k0-dz} & \NUM{st-oursx-best-avt-v-k0-pholm} & \NUM{st-oursx-best-avt-v-k0-tgtpos} & \NUM{st-oursx-best-avt-v-k5-mean} & \NUM{st-oursx-best-avt-v-k5-dz} & \NUM{st-oursx-best-avt-v-k5-pholm} & \NUM{st-oursx-best-avt-v-k5-tgtpos} \\
AVT & T & \NUM{st-oursx-best-avt-t-k0-mean} & \NUM{st-oursx-best-avt-t-k0-dz} & \NUM{st-oursx-best-avt-t-k0-pholm} & \NUM{st-oursx-best-avt-t-k0-tgtpos} & \NUM{st-oursx-best-avt-t-k5-mean} & \NUM{st-oursx-best-avt-t-k5-dz} & \NUM{st-oursx-best-avt-t-k5-pholm} & \NUM{st-oursx-best-avt-t-k5-tgtpos} \\
AVT & 融合 & \NUM{st-oursx-best-avt-f-k0-mean} & \NUM{st-oursx-best-avt-f-k0-dz} & \NUM{st-oursx-best-avt-f-k0-pholm} & \NUM{st-oursx-best-avt-f-k0-tgtpos} & \NUM{st-oursx-best-avt-f-k5-mean} & \NUM{st-oursx-best-avt-f-k5-dz} & $<$0.001 & \NUM{st-oursx-best-avt-f-k5-tgtpos} \\
AV & 融合 & \NUM{st-oursx-best-av-f-k0-mean} & \NUM{st-oursx-best-av-f-k0-dz} & \NUM{st-oursx-best-av-f-k0-pholm} & \NUM{st-oursx-best-av-f-k0-tgtpos} & \NUM{st-oursx-best-av-f-k5-mean} & \NUM{st-oursx-best-av-f-k5-dz} & \NUM{st-oursx-best-av-f-k5-pholm} & \NUM{st-oursx-best-av-f-k5-tgtpos} \\
VT & 融合 & \NUM{st-oursx-best-vt-f-k0-mean} & \NUM{st-oursx-best-vt-f-k0-dz} & $<$0.001 & \NUM{st-oursx-best-vt-f-k0-tgtpos} & \NUM{st-oursx-best-vt-f-k5-mean} & \NUM{st-oursx-best-vt-f-k5-dz} & $<$0.001 & \NUM{st-oursx-best-vt-f-k5-tgtpos} \\
AT & 融合 & \NUM{st-oursx-best-at-f-k0-mean} & \NUM{st-oursx-best-at-f-k0-dz} & $<$0.001 & \NUM{st-oursx-best-at-f-k0-tgtpos} & \NUM{st-oursx-best-at-f-k5-mean} & \NUM{st-oursx-best-at-f-k5-dz} & $<$0.001 & \NUM{st-oursx-best-at-f-k5-tgtpos} \\
AVT（第2） & A & \NUM{st-oursx-best-w2-a-k0-mean} & \NUM{st-oursx-best-w2-a-k0-dz} & \NUM{st-oursx-best-w2-a-k0-pholm} & \NUM{st-oursx-best-w2-a-k0-tgtpos} & \NUM{st-oursx-best-w2-a-k5-mean} & \NUM{st-oursx-best-w2-a-k5-dz} & $<$0.001 & \NUM{st-oursx-best-w2-a-k5-tgtpos} \\
AVT（第2） & V & \NUM{st-oursx-best-w2-v-k0-mean} & \NUM{st-oursx-best-w2-v-k0-dz} & \NUM{st-oursx-best-w2-v-k0-pholm} & \NUM{st-oursx-best-w2-v-k0-tgtpos} & \NUM{st-oursx-best-w2-v-k5-mean} & \NUM{st-oursx-best-w2-v-k5-dz} & \NUM{st-oursx-best-w2-v-k5-pholm} & \NUM{st-oursx-best-w2-v-k5-tgtpos} \\
AVT（第2） & T & \NUM{st-oursx-best-w2-t-k0-mean} & \NUM{st-oursx-best-w2-t-k0-dz} & \NUM{st-oursx-best-w2-t-k0-pholm} & \NUM{st-oursx-best-w2-t-k0-tgtpos} & \NUM{st-oursx-best-w2-t-k5-mean} & \NUM{st-oursx-best-w2-t-k5-dz} & $<$0.001 & \NUM{st-oursx-best-w2-t-k5-tgtpos} \\
AVT（第2） & 融合 & \NUM{st-oursx-best-w2-f-k0-mean} & \NUM{st-oursx-best-w2-f-k0-dz} & \NUM{st-oursx-best-w2-f-k0-pholm} & \NUM{st-oursx-best-w2-f-k0-tgtpos} & \NUM{st-oursx-best-w2-f-k5-mean} & \NUM{st-oursx-best-w2-f-k5-dz} & $<$0.001 & \NUM{st-oursx-best-w2-f-k5-tgtpos} \\
\bottomrule
\end{tabular}
\end{table}


\bibliographystyle{../IEEEtran}
\bibliography{../references}
\end{document}
